% !TeX root = ../../relatorio-projeto-aplicado-exemplo-undf.tex
\section*{Contexto e objetivos}
P01 valida solicitações isoladas de reserva de laboratório: campos vazios,
horários incoerentes e participantes acima da capacidade. O objetivo foi
implementar regras e conferir seus retornos com entradas sintéticas.
Laboratório e participantes são fictícios; os artefatos são executáveis.
O registro apoia a escrita do portfólio \citeind{manualundf}.

\section*{O que realizei e como}
No cenário, Estudante A definiu regras, implementou o programa e executou os
casos; Estudante B revisou regras e documentação. Esses papéis não atribuem
a pessoas reais a produção dos arquivos. O programa interpreta horários e
retorna o primeiro impedimento ou a aceitação pelas regras implementadas.

\section*{Decisões, dificuldades e aprendizados}
O escopo foi uma solicitação por vez, sem interface, armazenamento ou
conflitos entre reservas. O protótipo não é uma agenda em funcionamento.
É preciso relacionar regra, entrada e retorno; poucos casos não demonstram
ausência de defeitos nem adequação ao uso institucional.

\section*{Entregas, resultados e limites}
Foram produzidos o programa e seis casos sintéticos. A execução confirmou
dois retornos de aceitação e quatro recusas por motivos distintos.
\begin{undftabela}[H]{Casos registrados em P01}{casos-registro}[Elaboração própria a partir dos artefatos didáticos]
\begin{tabular}{@{}p{0.08\linewidth}p{0.46\linewidth}p{0.36\linewidth}@{}}
\hline
\textbf{Caso} & \textbf{Entrada examinada} & \textbf{Retorno observado} \\
\hline
C1 & 12 participantes para 20 vagas & Aceita \\
C2 & Solicitante vazio & Campos obrigatórios ausentes \\
C3 & Início às 25:00 & Horário inválido \\
C4 & Fim anterior ao início & Intervalo inválido \\
C5 & 21 participantes para 20 vagas & Capacidade excedida \\
C6 & 20 participantes para 20 vagas & Aceita no limite \\
\hline
\end{tabular}
\end{undftabela}
Não houve pesquisa com pessoas, implantação em laboratório ou medição de
impacto. A aceitação significa apenas passar pelas regras; não confirma
reserva, disponibilidade da sala ou autorização de acesso.

\section*{Evidências para recuperar depois}
O GitHub indicado é o acervo deste exemplo, não de uma equipe real.
O commit \texttt{685e69e} identifica os artefatos; os caminhos abaixo partem
da raiz do acervo. Para seu projeto, indique repositório próprio com README
e condições de acesso documentadas.

P01-E01 é o arquivo \texttt{validar\_reserva.py}; P01-E02 é \texttt{casos.json}.
Ambos estão em \texttt{exemplo/projeto-aplicado/artefatos/}. P01-E03 é a saída
reproduzível da execução, que informa seis casos com resultado esperado.
Na raiz do repositório, com Python 3 disponível, execute:
\begin{quote}\small\ttfamily
python exemplo/projeto-aplicado/artefatos/validar\_reserva.py
\end{quote}
Os arquivos não contêm dados de pessoas reais. O programa lê os casos sem
alterá-los e interrompe a execução se houver divergência.

\section*{Dúvidas e pontos a retomar na escrita}
Ao consolidar o portfólio, explicar o alcance de seis casos e a diferença
entre teste técnico e avaliação com usuários. Relacionar o limite de escopo
à contribuição individual e registrar que as funções excluídas não foram
entregues. Conservar esses limites ao comparar P01 com outros projetos.
